% ADR: Disk-Backed Snapshot Storage for Long Simulations
% Source: docs/adr-disk-storage.md (migrated to LaTeX)
% Uses macros from preamble.tex

\section{ADR: Disk-Backed Snapshot Storage for Long Simulations}
\label{sec:adr-disk-storage}

\textbf{Status:} Accepted \quad
\textbf{Date:} 2026-02-16

\textbf{Note:} This ADR predates the SUNDIALS migration (Feb 2026). References to py-pde reflect the architecture at time of writing; the solver layer now uses SUNDIALS IDA/CVODE + native numpy operators.

\textbf{Context:} TIDAL simulation pipeline stores entire time history in RAM, killing long-running simulations.


\subsection{Problem Statement}
\label{sec:adr-problem}

TIDAL stores \textbf{every snapshot in RAM} twice:

\begin{enumerate}
  \item \textbf{During simulation}: py-pde's \texttt{MemoryStorage} accumulates all snapshots in a Python list.
  \item \textbf{During measurement}: \texttt{SimulationData} eagerly loads the full NPZ into contiguous arrays.
\end{enumerate}

Memory formula: \texttt{n\_snapshots $\times$ (n\_fields + n\_momenta) $\times$ prod(grid\_shape) $\times$ 8 bytes}

\begin{table}[htbp]
\centering
\caption{Memory usage for representative simulations.}
\label{tab:adr-memory}
\begin{tabular}{@{}llccc@{}}
\toprule
Example & Grid & Snapshots & Slots & Memory \\
\midrule
coupled\_proca (2+1D) & $96^2$ & 1{,}250 & 10 & \textbf{880 MB} \\
gravitational\_waves (3+1D) & $64^3$ & 500 & 14 & \textbf{14 GB} \\
\bottomrule
\end{tabular}
\end{table}

At 3+1D grid sizes or long time series, the process is OOM-killed well before the simulation completes.


\subsection{Requirements}
\label{sec:adr-requirements}

\begin{itemize}
  \item \textbf{O(1) memory} in snapshot count for both writing and reading.
  \item \textbf{Framework-agnostic}: format must be readable from Python, Julia, C, Rust (future solver migration).
  \item \textbf{HPC-ready}: works on shared filesystems (NFS, Lustre, GPFS), crash-resilient (HPC wall-time kills).
  \item \textbf{Zero new dependencies}: only numpy.
  \item \textbf{Backward compatible}: existing NPZ workflows must still work.
\end{itemize}


\subsection{Alternatives Evaluated}
\label{sec:adr-alternatives}

\subsubsection{HDF5 (PDEBench Approach)}

PDEBench uses HDF5 for PDE simulation datasets with chunked access and per-dataset compression.

\textbf{Pros:} Mature, multi-language (C/Fortran/Julia/Python), chunked random access, widely used in HPC.

\textbf{Cons:} Requires \texttt{h5py} dependency (C library build), breaks zero-new-deps principle. File locking can be problematic on some shared filesystems.

\textbf{Decision:} Rejected. Adds non-trivial dependency; HDF5 file locking has known issues on NFS.

\subsubsection{Zarr}

Zarr offers chunked + compressed N-dimensional array storage with cloud support (S3, GCS).

\textbf{Pros:} Modern Python API, excellent multi-axis slicing, compression (Blosc, zstd), cloud-friendly, append-friendly.

\textbf{Cons:} Adds \texttt{zarr} + \texttt{numcodecs} dependencies, Python-centric (no native Julia/C reader), compression provides poor ratio on float64 physics data (typically 1.5--2$\times$ vs 8$\times$ for integer data).

\textbf{Decision:} Deferred as future extension. If cloud storage or compression becomes needed, a \texttt{ZarrWriter} can implement the same \texttt{append()}/\texttt{close()} API without changing measurement code.

\subsubsection{py-pde FileStorage}

py-pde provides built-in \texttt{FileStorage} for HDF5-based streaming during simulation.

\textbf{Pros:} Direct integration with solver, streaming.

\textbf{Cons:} Still needs \texttt{h5py}, couples storage format to py-pde (blocks Julia migration), py-pde-specific API.

\textbf{Decision:} Rejected. Framework coupling is a non-starter for planned Julia migration.

\subsubsection{Incremental Measurement During Simulation}

Compute energy/conversion at each snapshot, discard field data immediately. $O(1)$ memory.

\textbf{Pros:} Absolute minimum memory.

\textbf{Cons:} User cannot do post-hoc analysis (new measurement types, debugging, visualization). Defeats the purpose of \texttt{tidal measure} CLI.

\textbf{Decision:} Rejected. Full field data must be available for flexible post-hoc analysis.

\subsubsection{NPZ with Memory Mapping}

numpy supports \texttt{mmap\_mode} when loading NPZ files.

\textbf{Pros:} No format change needed.

\textbf{Cons:} Individual arrays within a ZIP-based NPZ are not contiguous on disk. Memory mapping through the ZIP layer is unreliable for large files. The per-snapshot keying (\texttt{phi\_0\_t0}, \texttt{phi\_0\_t1}, \ldots) prevents true contiguous time-axis access.

\textbf{Decision:} Rejected. NPZ's ZIP container prevents efficient mmap.

\subsubsection{Directory of \texttt{.npy} Files with Memory Mapping (Chosen)}

Pre-allocate one \texttt{.npy} file per field with shape \texttt{(n\_snapshots, *grid\_shape)}, write via \texttt{numpy.memmap(mode='w+')}, read via \texttt{numpy.load(mmap\_mode='r')}.

\paragraph{Pros.}
\begin{itemize}
  \item Zero dependencies (pure numpy).
  \item C-contiguous time-first layout matches our sequential access pattern (snapshot-by-snapshot iteration).
  \item Julia reads \texttt{.npy} natively (\texttt{NPZ.jl}, \texttt{Mmap.mmap()}). C/Rust have \texttt{.npy} parsers.
  \item Pre-allocated files provide crash resilience: all flushed snapshots survive HPC wall-time kills.
  \item Exact snapshot count is known at start (\texttt{int(t\_end / interval) + 1}), so no upper-bound estimation or truncation needed.
  \item Works on any POSIX filesystem, Windows, VMs, HPC clusters.
\end{itemize}

\paragraph{Cons.}
\begin{itemize}
  \item No compression (float64 data compresses poorly anyway; ${\sim}1.5$--$2\times$ with Blosc).
  \item No random-access chunking (but our access is sequential along time axis, so this is fine).
  \item One file per field (manageable; typical simulations have 4--10 fields).
\end{itemize}

\textbf{Decision:} Accepted.


\subsection{Design}
\label{sec:adr-design}

\subsubsection{Directory Layout}

\begin{lstlisting}[style=bash]
output_dir/
  metadata.json          # Grid, parameters, field list, snapshot count
  times.npy              # shape (n_snapshots,), float64
  phi_0.npy              # shape (n_snapshots, *grid_shape), float64
  chi_0.npy              # shape (n_snapshots, *grid_shape), float64
  pi_phi_0.npy           # shape (n_snapshots, *grid_shape), float64
  pi_chi_0.npy           # shape (n_snapshots, *grid_shape), float64
\end{lstlisting}

\subsubsection{metadata.json Schema}

\begin{lstlisting}[style=json]
{
  "version": 1,
  "n_snapshots": 1250,
  "grid_spacing": [0.5208, 0.5208],
  "grid_bounds": [
    [0, 50],
    [0, 50]
  ],
  "grid_shape": [96, 96],
  "periodic": [true, true],
  "parameters": { "mA": 1.0, "mB": 1.0, "g": 0.1 },
  "spec_path": "examples/data/coupled_proca.json",
  "fields": ["A_1", "A_2", "B_1", "B_2"],
  "momenta": ["A_1", "A_2", "B_1", "B_2"],
  "dtype": "float64"
}
\end{lstlisting}

\subsubsection{Exact Snapshot Count}

Since \texttt{t\_end} and \texttt{snapshot\_interval} are both known at simulation start:

\begin{lstlisting}[style=python]
n_snapshots = int(t_end / snapshot_interval) + 1   # +1 for initial state at t=0
\end{lstlisting}

This is exact. No upper-bound estimation, no truncation needed. The \texttt{.npy} files are pre-allocated with the correct shape.

\subsubsection{Writing: SnapshotWriter}

\texttt{SnapshotWriter} uses \texttt{numpy.lib.format.open\_memmap()} to pre-allocate proper \texttt{.npy} files (with headers), then writes each snapshot in-place:

\begin{lstlisting}[style=python]
# Each append() is O(grid_size) -- writes to the next row in the memmap
writer = SnapshotWriter(output_dir, field_names, momentum_names,
                         grid_shape, n_snapshots, ...)
for t, fields, momenta in simulation:
    writer.append(t, fields, momenta)  # O(1) memory
writer.close()  # writes metadata.json
\end{lstlisting}

Key properties:

\begin{itemize}
  \item Each \texttt{append()} writes to pre-allocated memmap rows. $O(1)$ memory, $O(\text{grid\_size})$ I/O.
  \item \texttt{flush()} on each append for crash resilience.
  \item \texttt{metadata.json} written at \texttt{close()}. If missing (crash), snapshot count is recoverable from \texttt{times.npy} by finding the last non-zero entry.
\end{itemize}

\subsubsection{Reading: SimulationData.from\_directory()}

\begin{lstlisting}[style=python]
data = SimulationData.from_directory(output_dir, spec)
# data.fields["phi_0"] is a np.memmap -- O(1) RAM
# Only pages actually accessed are loaded by the OS page cache
\end{lstlisting}

\texttt{mmap\_mode="r"} gives zero-copy, read-only access. The OS manages which pages are in RAM via the page cache. Sequential access (iterate snapshots) is optimal because the data is C-contiguous along the time axis.

\subsubsection{CLI Integration}

\begin{itemize}
  \item \texttt{tidal simulate spec.json --output /tmp/results/} creates a snapshot directory (disk-backed, $O(1)$ memory).
  \item \texttt{tidal measure /tmp/results/ --spec spec.json} loads from the snapshot directory.
  \item \texttt{SimulationData.load(path, spec)} is the universal entry point (directory only; NPZ support was removed).
\end{itemize}

\subsubsection{Crash Recovery}

When \texttt{metadata.json} is missing (writer wasn't closed due to crash/kill):

\begin{enumerate}
  \item Look for \texttt{times.npy} in the directory.
  \item Memory-map it read-only.
  \item Find the last index where \texttt{times[i] > 0} (unwritten entries are 0.0 from memmap pre-allocation).
  \item Use that count as \texttt{n\_snapshots}. All prior snapshots are intact.
\end{enumerate}

\subsubsection{Access Pattern Analysis}

\begin{table}[htbp]
\centering
\caption{Measurement access patterns and mmap behaviour.}
\label{tab:adr-access}
\begin{tabular}{@{}lll@{}}
\toprule
Measurement Function & Access Pattern & Mmap Behaviour \\
\midrule
\texttt{compute\_energy\_timeseries} & Sequential over time & Optimal: sequential page faults \\
\texttt{compute\_conversion\_probability} & Sequential over time & Optimal: same as above \\
\texttt{compute\_spectrum} & All snapshots of one field & Full scan: pages loaded sequentially \\
\texttt{compute\_dispersion} & Full field + temporal FFT & Full scan: all pages loaded \\
\texttt{compute\_mixing\_length} & Derived from conversion result & No direct field access \\
\bottomrule
\end{tabular}
\end{table}

All access patterns are sequential along the first (time) axis, which is exactly what C-contiguous memmap optimizes for.


\subsection{HPC and Cross-Platform Compatibility}
\label{sec:adr-hpc}

\begin{itemize}
  \item \textbf{No dependencies beyond numpy}: Works on any system with Python + numpy.
  \item \textbf{POSIX + Windows}: \texttt{numpy.memmap} abstracts OS-level mmap calls.
  \item \textbf{Shared filesystems} (NFS, Lustre, GPFS): Sequential I/O pattern (write forward, read forward) is optimal. No random access, no file locking.
  \item \textbf{Crash resilience}: Files pre-allocated at creation. \texttt{append()} flushes per snapshot. If job is killed (HPC wall-time), all flushed snapshots are intact.
  \item \textbf{Julia interop}: Julia's \texttt{NPZ.jl} or \texttt{Mmap.mmap()} reads \texttt{.npy} files natively. A Julia solver can write the same directory layout.
  \item \textbf{Large files}: 64-bit systems handle files ${>}2\,\text{GB}$. A 3+1D $64^3 \times 500$ snapshot field $= 1\,\text{GB}$ per \texttt{.npy} file, well within limits.
  \item \textbf{Virtual machines}: Standard filesystem operations. No special kernel modules or shared memory needed.
\end{itemize}


\subsection{Future Extensions}
\label{sec:adr-future}

\begin{enumerate}
  \item \textbf{Zarr backend}: If compression or cloud storage is needed, a \texttt{ZarrWriter} can implement the same \texttt{append()}/\texttt{close()} API. Measurement code is unchanged because \texttt{SimulationData} provides the abstraction layer.
  \item \textbf{Constraint field exclusion}: \texttt{SnapshotWriter} can skip constraint fields (\texttt{time\_derivative\_order == 0}) that are re-derivable from dynamical fields, saving ${\sim}20\%$ storage for systems like coupled\_proca.
  \item \textbf{Parallel writing}: For domain-decomposed simulations, each rank writes its own subdirectory; a merge step concatenates along spatial axes.
\end{enumerate}


\subsection{References}
\label{sec:adr-references}

\begin{itemize}
  \item PDEBench data format (HDF5) --- similar PDE simulation storage problem, chose HDF5.
  \item mmap vs Zarr/HDF5 performance --- mmap is optimal when access aligns with storage layout.
  \item \texttt{numpy.lib.format.open\_memmap} --- creates proper \texttt{.npy} with header.
  \item py-pde \texttt{CallbackTracker} --- \texttt{CallbackTracker(func, interrupts=interval)} for streaming snapshots.
\end{itemize}
